\chapter{State of the Art}

\section{Information retrieval and documentary indexing}

The intended outcome of Index\_Insight is not simply a collection of text
files. It is a corpus that can be explored by theme, initiative type,
territorial scale, holder, and source. This objective places the project at the
intersection of information retrieval and documentary indexing: retrieval makes
documents findable, while indexing makes their relevant properties available
for filtering and comparison.

Lexical retrieval remains an appropriate starting point for this type of
collection. TF--IDF weights terms according to their importance in a document
and their rarity in the corpus \citep{salton1988term}; BM25 extends this idea
with a probabilistic ranking model that accounts for term frequency and document
length \citep{robertson2009probabilistic}. Both approaches are transparent,
efficient, and well suited to a future keyword search interface. They are not,
however, sufficient on their own when two documents describe a similar
initiative with different vocabulary, or when a user needs to filter documents
by information that is only stated in the body text.

Dense retrieval represents queries and documents in a shared vector space and
can retrieve semantically related material even when their wording differs
\citep{karpukhin2020dense}. This is a relevant future direction for INSIGHT,
especially for exploratory queries formulated in everyday language. It is not
implemented in the current prototype. The immediate requirement is instead to
produce reliable structured metadata from source documents; those metadata can
later support faceted search, lexical retrieval, or a hybrid lexical--semantic
interface.

Metadata are central to this distinction. Titles, dates, organisations,
territories, target audiences, and themes make documentary collections easier
to navigate than full text alone. In the present corpus, these elements are not
consistently available as webpage metadata. They must be recovered from HTML
pages, PDF-derived text, and action sheets, then retained with a link to their
source. The quality of this metadata extraction is therefore more important to
the current project than the choice of a final ranking algorithm.

\section{Structured extraction from heterogeneous documents}

Information extraction from the present corpus is difficult because the same
piece of information takes different forms across sources. A territorial scale
may be stated in a heading, embedded in a sentence, or only implied by a place
name. A date can refer to the publication of a report, the start of an
initiative, or a future event. Websites also vary in HTML structure, while PDF
conversion may introduce fragmented sentences, repeated headers, and encoding
artefacts. Source-specific collection is therefore a practical necessity rather
than an implementation detail \citep{mitchell2018web}.

Rule-based extraction is a natural baseline in this setting. Regular
expressions, headings, lexicons, and keyword rules are transparent: a predicted
value can be traced to a matching expression and the rule can be corrected.
Such methods are effective when documents follow stable templates. Their main
limitation is maintenance. A rule designed for one editorial structure can fail
when the information appears elsewhere or is expressed differently
\citep{sarawagi2008information,grishman2019information}. This limitation is
particularly relevant here because the corpus combines six source collections
with substantially different document formats.

Preprocessing reduces some of this variability before extraction. In the
pipeline, it includes text cleaning, whitespace and file-name normalisation,
and the separation of raw from cleaned documents. More general operations such
as tokenisation, lemmatisation, and normalisation are standard components of
NLP pipelines because they make textual units easier to compare and process
\citep{manning1999foundations}. They cannot, however, resolve whether an
organisation mentioned in a document is the initiative holder or merely a
partner. That decision requires contextual interpretation.

Structured extraction also requires a schema. A document can be represented by
fields such as target audience, territorial scale, initiative type, territory,
holder, date, and theme. Controlled vocabularies are useful where comparable
values are needed: they reduce spelling variation and make filtering possible
\citep{hedden2016taxonomy}. They should not erase the original wording. For
this reason, Index\_Insight preserves the extracted values and adds a separate,
deterministic normalisation layer for themes and initiative types.

\section{LLM-assisted extraction and constrained generation}

Transformer models make it possible to use the context surrounding a candidate
value rather than relying exclusively on an isolated keyword. The transformer
architecture introduced by \citet{vaswani2017attention}, and contextual models
such as BERT \citep{devlin2019bert}, established the importance of contextual
representations for many NLP tasks. Large language models extend this approach:
a task can be specified through an instruction and a desired output format
rather than through a task-specific classifier or a large set of source-specific
rules \citep{brown2020language}.

For Index\_Insight, the relevant use of an LLM is structured information
extraction, not open-ended text generation. The model receives a French prompt
that specifies a JSON schema, the expected type of each field, controlled values
for target audience and territorial scale, and an explicit instruction not to
invent unsupported information. Prompting and output constraints are important
because they make generated records easier to parse, compare, and review
\citep{liu2023pre}. They also allow empty values to be used when the document
does not support a reliable answer.

Constrained output does not make an LLM output factual by default. A syntactically
valid JSON record may still contain an unsupported interpretation, omit a
relevant value, or select an overly broad category. This concern is central in
institutional and research contexts, where fluent language can conceal
uncertainty \citep{bommasani2021opportunities}. The project consequently
retains raw model responses and source provenance, applies deterministic
normalisation after extraction, and evaluates selected fields against human
annotations. The LLM is thus treated as a component of a traceable workflow,
not as an autonomous source of knowledge.

\section{Comparable systems}

Existing systems address parts of the problem tackled by Index\_Insight, but
they usually assume either that metadata already exist or that documents belong
to a relatively uniform collection.

Documentary catalogues and digital libraries organise resources through
bibliographic or descriptive metadata and offer faceted navigation by author,
date, subject, institution, or place. Their strength is the use of stable
schemas and controlled vocabularies. They do not, by themselves, extract the
metadata needed for this project from long web pages, reports, and action
sheets.

Search systems based on inverted indexes, TF--IDF, or BM25 can efficiently
retrieve documents containing query terms \citep{salton1988term,robertson2009probabilistic}.
Dense retrieval can complement them when query and document vocabulary differ
\citep{karpukhin2020dense}. Neither approach creates the territorial,
organisational, temporal, and thematic fields required for the comparative
exploration envisaged by INSIGHT.

Rule-based information-extraction pipelines use templates, headings, regular
expressions, and lexicons to identify candidate values. They offer useful
traceability but require continuous adaptation when formats change
\citep{sarawagi2008information,grishman2019information}. Prompt-guided LLM
systems instead request a predefined structure directly from the document. They
are more flexible in the face of heterogeneous phrasing, but their outputs must
be constrained and validated \citep{liu2023pre}.

\begin{table}[H]
\centering
\begin{tabular}{p{3.2cm}p{3.5cm}p{3.5cm}p{3.4cm}}
\toprule
Approach & Main strength & Main limitation for this corpus & Position of Index\_Insight \\
\midrule
Documentary catalogue & Stable metadata and faceted navigation & Assumes metadata already exist or are manually curated & Reuses the principle of controlled metadata, but automates part of their production \\

Lexical or dense retrieval & Efficient retrieval and exploration of large collections & Does not itself extract reliable documentary metadata & Future use of the structured corpus; not yet a deployed search engine \\

Rule-based extraction & Transparent and inspectable decisions & Fragile across heterogeneous document structures & Baseline and quality-control layer of the pipeline \\

LLM-assisted extraction & Context-sensitive structuring of varied documents & Risk of unsupported inference and output inconsistency & Constrained local extraction, evaluated against human reference annotations \\
\bottomrule
\end{tabular}
\caption{Positioning of Index\_Insight relative to comparable system families.}
\label{tab:comparable-systems}
\end{table}

\section{Synthesis of the state of the art}

The literature supports a hybrid design for a corpus of this kind. Rule-based
methods offer traceability and remain useful for explicit lexical cues, but they
are difficult to generalise across heterogeneous layouts. LLM-assisted
extraction can use contextual information to select and format metadata, but it
requires constraints, provenance, and empirical evaluation. Controlled
vocabularies make the resulting records comparable, while future retrieval
components can use both full text and structured fields.

Index\_Insight adopts this division of labour. The pipeline preserves raw
sources and intermediate rule-based outputs, uses a local LLM to create JSON
records, enriches them with provenance, and normalises selected fields without
overwriting the original outputs. This design does not assume that a completed
field is correct. It distinguishes technical completion from factual accuracy
and makes human validation part of the workflow. The following chapters detail
how this design was implemented and evaluated on the collected corpus.
